\chapter*{Introduction générale}
\phantomsection
\addcontentsline{toc}{chapter}{Introduction générale}
\markboth{Introduction générale}{Introduction générale}

Dans une activité de boulangerie-restauration, la décision de produire ne peut
pas être reportée au moment où la demande se manifeste. Les équipes doivent
préparer plusieurs références avant l'ouverture ou plusieurs heures avant leur
vente, alors que la demande varie selon le jour de la semaine, la saison, les
fêtes, les événements et les commandes exceptionnelles. Une quantité trop faible
provoque des ruptures et une perte de chiffre d'affaires ; une quantité trop
élevée augmente les invendus, immobilise des matières premières et dégrade la
rentabilité. Cette tension rend la prévision directement opérationnelle : elle
doit être compréhensible, actualisable et reliée aux recettes de production.

Le contexte étudié est celui d'un point de vente PAUL au Maroc. Les informations
historiques existent sous plusieurs formes : exports mensuels, états journaliers
issus du logiciel de caisse, recettes techniques, prix de matières et événements
connus. Avant le projet, ces éléments ne formaient pas une chaîne décisionnelle
unique. L'historique permettait d'observer les ventes passées, mais il ne répondait
pas automatiquement aux questions quotidiennes : combien produire demain,
quelles références sont incertaines, quelles matières commander, quel effet
attendre d'une fête ou d'une commande client, et comment contrôler la qualité
des prévisions ?

La problématique peut donc être formulée ainsi : \emph{comment concevoir un
système local qui transforme des données de vente hétérogènes en prévisions
journalières et mensuelles fiables, puis en recommandations de production et
d'approvisionnement traçables, tout en restant utilisable par une équipe non
spécialiste de la science des données ?}

Pour répondre à cette problématique, nous avons développé \NomProjet. La solution
réunit une chaîne de préparation de données, plusieurs méthodes de prévision, une
validation temporelle sans fuite d'information, des ajustements métier, un calcul
MRP à partir des nomenclatures, des exports opérationnels et un tableau de bord
web. Le système est conçu pour fonctionner localement afin de préserver la
confidentialité des ventes. Il peut également se mettre à jour automatiquement
depuis la base SQL de la caisse lorsque l'infrastructure de production est
disponible.

Ce rapport est organisé en quatre chapitres. Le premier présente le cadre du
projet, l'existant, ses limites, la solution retenue et la démarche de travail. Le
deuxième formalise les acteurs, les besoins fonctionnels et non fonctionnels ainsi
que les principaux cas d'utilisation. Le troisième expose la conception dynamique,
l'architecture logicielle, le modèle de données et les choix de modélisation. Le
quatrième décrit la réalisation, les interfaces principales, le déploiement, les
tests et les résultats mesurés. Une conclusion générale synthétise enfin les
apports, les limites actuelles et les perspectives d'évolution.
